\chapter{DTFT}

\begin{figure}[h]
    \centering
    \includegraphics[scale = 0.115]{Discrete Time Fourier Transform (DTFT) explained visually 2-22 screenshot.png}
\end{figure}  

\newpage 

\section{Risposta in frequenza dei sistemi LTI}

\footnote{ 
    Slide del prof | 3 - DTFT | pag 2 - 4\\
    Appunti mancanti | 2025-10-07 
}

Dato: 

\begin{itemize}
    \item un sistema LTI (Linear Time-Invariant, cioè un sistema lineare che non cambia le sue caratteristiche nel tempo) 
    \item una sequenza di ingresso $x \left[ n\right]$ una sinusoide complessa (in particolare un'esponenziale complesso puro) di pulsazione $\omega$, quindi $e^{\jmath \omega n}$ 
\end{itemize}

allora, l'uscita del sistema LTI sarà ancora una sinusoide complessa di pulsazione $\omega$. \newline 

Dal punto di vista visuale avremo il seguente schema: 

\begin{figure}[h]
    \centering
    \includegraphics[scale = 0.4]{Risposta sinusoidale di un sistema LTI.png}
\end{figure}  

\newpage 

\begin{tcolorbox}
    Per avere un ripasso al volo perchè anche nel DSP, come anche a Teoria dei Segnali e nel corso di Controlli Automatici si fanno molto uso degli esponenziali, 
    vi lascio questo video: \newline 

    The Physics of Euler's Formula | Laplace Transform Prelude by 3Blue1Brown \\
    \url{https://www.youtube.com/watch?v=-j8PzkZ70Lg&t=111s} \newline 

    In breve, si utilizza il piano dei complessi come il seguente: \newline 

    \includegraphics[scale = 0.15]{The Physics of Euler's Formula _ Laplace Transform Prelude 5-36 screenshot.png}
    
    per rappresentare delle sinusoidi nel tempo: \newline

    \includegraphics[scale = 0.5]{sinusoide nel tempo.jpeg}
    
    Dal corso di Teoria dei segnali, sappiamo che la pulsazione $\omega$ vale: 

    {
        \Large 
        \begin{equation}
            \omega = 2 \pi f
        \end{equation}
    }

    dove f è la frequenza. \newline 

    La frequenza f, nel piano dei complessi, indica quanti "giri" fa il vettore intorno al cerchio unitario al secondo. 
\end{tcolorbox}

\newpage 

Di seguito la dimostrazione formale dell'uscita: 

{
    \Large 
    \begin{equation}
        \setlength{\jot}{10pt} % affecting the line spacing in the environment
        \begin{split}
            y \left[ n\right] 
            &= 
            \sum_{k = - \infty}^{+ \infty}
            h \left[ k\right]
            \cdot 
            e^{\jmath \omega \left( n - k\right)}
            \\
            &=
            \sum_{k = - \infty}^{+ \infty}
            h \left[ k\right]
            \cdot 
            e^{\jmath \omega n}
            \cdot 
            e^{-\jmath \omega k}
            \\
            &=
            \left\{
            \sum_{k = - \infty}^{+ \infty}
            h \left[ k\right]
            \cdot 
            e^{-\jmath \omega k}
            \right\}
            \cdot 
            e^{\jmath \omega n}
            \\
            &= 
            H\left( e^{\jmath \omega}\right)
            \cdot 
            e^{\jmath \omega n}
        \end{split}
    \end{equation}
}

Il termine: 

{
    \Large 
    \begin{equation}
        H\left( e^{\jmath \omega}\right)
        = 
        \sum_{k = - \infty}^{+ \infty}
        h \left[ k\right]
        \cdot 
        e^{-\jmath \omega k}    
    \end{equation}
} 

prende il nome di risposta in frequenza, la quale è periodica ogni multiplo di $2 \pi$. \newline 

Quindi, la riposta in frequenza $H\left( e^{\jmath \omega}\right)$ è una funzione complessa (perchè la si può rappresentare nel piano dei complessi) 
che dipende dalla pulsazione $\omega$. \newline 

In particolare la risposta in frequenza $H\left( e^{\jmath \omega}\right)$, siccome la si può rappresentare sul piano dei complessi il quale a sua volta è formato da un asse reale e da un asse immaginario, 
si può dividere in: 

\begin{itemize}
    \item riposta in ampiezza (in inglese prende il nome di "Magnitude") e lo si denota con il seguente simbolo $\abs{H\left( e^{\jmath \omega}\right)}$ 
    \item risposta in fase e lo si denota con il seguente simbolo $\sphericalangle H\left( e^{\jmath \omega}\right)$
\end{itemize}

\begin{tcolorbox}
    Nel piano dei fasori, ampiezza e fase sono rappresentati in questa maniera: \newline 

    \includegraphics[scale = 0.5]{complex plane magnitude and phase.png}
\end{tcolorbox}

\newpage 

\section{Trasformata di Fourier di Segnali Discreti} \label{sec: Trasformata di Fourier di Segnali Discreti}

\footnote{ 
    Slide del prof | 3 - DTFT | pag 5\\
    Appunti mancanti | 2025-10-07 \\ 
    Latex | 2025-10-13 Sera | pag 65
} 

Come studiato nel corso di Teoria dei segnali, 
un segnale analogico, che è continuo sia nelle ampiezze f(t) che nel tempo t, 
è possibile rappresentarlo anche nel dominio di Fourier con la seguente formula: 

{
    \Large 
    \begin{equation}
        F\left( \jmath \omega \right) 
        = 
        \int_{- \infty}^{+ \infty}
        f(t) \cdot e^{- \jmath \omega t} 
        dt 
    \end{equation}
} 

che prende il nome di Trasformata di Fourier. \newline 

La formula inversa, cioè l'anti-trasformata di Fourier, 
che è quella formula che riporta il segnale $F \left(\jmath \omega\right)$ dal dominio di Fourier al dominio nel tempo t, 
è la seguente: 

{
    \Large 
    \begin{equation}
        f \left( t\right) 
        = 
        \dfrac{1}{2 \pi}
        \int_{- \infty}^{+ \infty}
        F \left(\jmath \omega \right) \cdot e^{\jmath \omega t} 
        d\omega
    \end{equation}
}

Invece, nel dominio dei discreti, non useremo la trasformata di Fourier, bensì useremo questa formula: 

{
    \Large 
    \begin{equation}
        X \left( e^{\jmath \omega}\right)
        = 
        \sum_{n = - \infty}^{+ \infty}
        x \left[ n\right]
        \cdot 
        e^{- \jmath \omega n}
    \end{equation}
}

che prende il nome di DTFT (cioè Discrete-Time Fourier Transform). \newline 

Invece la formula inversa della DTFT (l'anti-DTFT) è la seguente: 

{
    \Large 
    \begin{equation}
        x \left[ n\right]
        = 
        \dfrac{1}{2 \pi} 
        \int_{- \pi}^{+ \pi}
        X \left(e^{\jmath \omega}\right)
        \cdot 
        e^{\jmath \omega n}
    \end{equation}
}

\begin{tcolorbox}
Nella formula dell'anti-trasformata della DTFT compare un integrale perchè, 
come dice il nome stesso, nella DTFT (cioè Discrete-Time Fourier Transform) 
il tempo è discretizzato mentre i valori nelle ampiezze sono continui. 
\end{tcolorbox}

$X \left(e^{\jmath \omega}\right)$ è periodica ogni multiplo di $2 \pi$. \newline 

Come la risposta in frequenza dei sistemi LTI, 
possiamo rappresentare $X \left(e^{\jmath \omega}\right)$ come: 

\begin{itemize}
    \item $\abs{X \left(e^{\jmath \omega}\right)}$, che è lo spettro di ampiezza di $x \left[ n\right]$ 
    \item $\sphericalangle X \left(e^{\jmath \omega}\right)$, che è lo spettro di fase di $x \left[ n\right]$
\end{itemize}

\newpage 

\subsection{Esistenza della trasformata} \label{subsec: Esistenza della trasformata}

\footnote{ 
    Slide del prof | 3 - DTFT | pag 6\\
    Appunti mancanti | 2025-10-07 \\ 
    Latex | 2025-10-13 Sera | pag 66 - 67
}

Come nel caso continuo studiato a Teoria dei Segnali, 
non tutte le sequenze $x \left[ n\right]$ sono trasformabili nel dominio di Fourier. \newline 

Di seguito alcune condizioni per cui $x \left[ n\right]$ è trasformabile in Fourier. \newline 

\textbf{Condizione di esistenza numero 1} \newline 

Se una sequenza $x \left[ n\right]$ è assolutamente sommabile, 
allora la sua trasformata di Fourier $X \left( e^{\jmath \omega }\right)$ esiste ed è funzione continua di $\omega$. \newline 

In formule: 

{
    \Large 
    \begin{equation}
        \sum_{n = - \infty}^{+ \infty}
        \abs{x \left[ n\right]} 
        < 
        + \infty
        \to 
        \text{ convergenza uniforme}
    \end{equation}
} 

che è una condizione sufficiente. \newline 

\begin{tcolorbox}
    Se non ti ricordi dall'esame di Analisi Matematica 1 cosa significa che una serie ha una convergenza uniforme: \newline 

    Convergenza uniforme by Wikipedia \\
    \url{https://it.wikipedia.org/wiki/Successione_di_funzioni#Convergenza_uniforme} 
\end{tcolorbox} 

\textbf{Condizione di esistenza numero 2} \newline 

Se una sequenza $x \left[ n\right]$ è quadraticamente sommabile, 
allora la sua trasformata di Fourier $X \left( e^{\jmath \omega }\right)$ esiste ed è funzione discontinua di $\omega$. \newline 

In formule: 

{
    \Large 
    \begin{equation} 
        \sum_{n = - \infty}^{+ \infty}
        \abs{x \left[ n\right]}^{2} 
        < 
        + \infty
        \to 
        \text{ convergenza non uniforme}
    \end{equation}
} 

che è una condizione sufficiente. \newline 

\begin{tcolorbox}
Per funzione discontinua in $\omega$ significa che la funzione avrà dei tratti continui e dei tranni non continui. \newline 

Ma, in questo corso, sarà più utile la prima condizione di esistenza che la seconda. 
\end{tcolorbox}

\newpage 

\textbf{Condizione di esistenza numero 3} \newline 

Se la sequenza $x \left[ n\right]$ non è ne assolutamente o quadraticamente sommabile, 
allora $X \left( e^{\jmath \omega}\right)$ può esistere in casi particolari, come ad esempio: 

{
    \Large 
    \begin{equation}
        \begin{split}
        x \left[ n \right]  
        &= 
        e^{\jmath \omega_0 n} 
        \text{ } 
        -\infty < n < + \infty
        \\
        &\updownarrow
        \\
        X \left( e^{\jmath \omega}\right)
        &=
        \sum_{k = - \infty}^{+ \infty}
        2 \pi \delta \left( \omega - \omega_0 + 2 \pi k\right)  
        \end{split}
    \end{equation}
}

dove, nella formula, $\delta$ indica un impulso unitario continuo (Dirac). \newline 

\begin{tcolorbox}
$\delta \left( \omega - \omega_0 + 2 \pi k\right) $ indica una distribuzione. 
\end{tcolorbox}
 
\newpage 

\section{Sistem LTI e Trasformata di Fourier}

\footnote{ 
    Slide del prof | 3 - DTFT | pag 7 - 8\\
    Appunti mancanti | 2025-10-07 
} 

La risposta in frequenza del sistema LTI è la trasformata di Fourier della risposta impulsiva: 

{
    \Large 
    \begin{equation}
        H\left( e^{\jmath \omega}\right)
        = 
        \sum_{k = - \infty}^{+ \infty}
        h \left[ k\right]
        \cdot 
        e^{-\jmath \omega k}    
    \end{equation}
} 

La relazione tra ingresso ed uscita in frequenza del sistema LTI è la seguente: 

{
    \Large 
    \begin{equation}
        Y \left( e^{\jmath \omega}\right) 
        = 
        H \left( e^{\jmath \omega}\right) 
        \cdot 
        X \left( e^{\jmath \omega}\right) 
    \end{equation}
}

(che è la stessa del caso continuo della trasformata di Fourier). \newline 

\newpage 

\subsection{Proprietà della trasformata di Fourier} \label{subsubsec: Proprietà della trasformata di Fourier}

\footnote{ 
    Slide del prof | 3 - DTFT | pag 9 - 10\\
    Appunti mancanti | 2025-10-07 \\ 
    Latex | 2025-10-13 Sera | pag 70
} 

Come nel caso continuo, anche nel caso discreto si applicano le proprietà della trasformata di Fourier. \newline 

Di seguito le proprietà principali. \newline 

\textbf{Linearità} \newline

Date: 

\begin{itemize}
    \item due sequenze di ingresso $x_1 \left[ n\right]$ e $x_2 \left[ n\right]$ 
    \item le rispettive rappresentazione in Fourier $X_1 \left(e^{\jmath \omega}\right)$ e $X_2 \left(e^{\jmath \omega}\right)$ 
    \item a e b come coefficienti moltiplicativi  
\end{itemize}

possiamo esprimere la proprietà di linearità: 

{
    \Large 
    \begin{equation}
        a \cdot x_1 \left[ n\right]
        + 
        b \cdot x_1 \left[ n\right]
        \leftrightarrow
        a \cdot X_1 \left(e^{\jmath \omega}\right) 
        +
        b \cdot X_1 \left(e^{\jmath \omega}\right)  
    \end{equation}
}

\textbf{Traslazione} \newline 

Se consideriamo $n_0$ un coefficiente di ritardo, 
possiamo esprimere la seguente proprietà di traslazione: 

{
    \Large 
    \begin{equation}
        x \left[ n - n_0\right]
        \leftrightarrow 
        X \left(e^{\jmath \omega}\right) \cdot e^{- \jmath \omega n_0}
    \end{equation}
}

Se invece consideriamo la pulsazione fondamentale $\omega_0$ , 
allora possiamo esprimere la seguente proprietà di traslazione in frequenza: 

{
    \Large 
    \begin{equation}
        x \left[ n\right] \cdot e^{\jmath \omega_0 n}
        \leftrightarrow 
        X \left(e^{\jmath \left(\omega - \omega_0\right)}\right)
    \end{equation}
}

\textbf{Rovesciamento nel tempo} \newline 

{
    \Large 
    \begin{equation}
        x \left[ -n\right]
        \leftrightarrow
        X \left(e^{- \jmath \omega}\right)
    \end{equation}
}

oppure 

{
    \Large 
    \begin{equation}
        x \left[ -n\right]
        \leftrightarrow
        X^{*} \left(e^{\jmath \omega}\right) 
        \text{ se }
        x \left[ n\right] \in \mathbb{R}
    \end{equation}
}

dove con $X^{*} \left(e^{\jmath \omega}\right)$ si indica il coniugato di $X \left(e^{- \jmath \omega}\right)$. \newline 

\textbf{Convoluzione} \newline 

Sapendo che, la formula di convoluzione è la seguente: 

{
    \Large 
    \begin{equation}
        x \left[ n\right]
        * 
        y \left[ n\right] 
        = 
        \sum_{k = -\infty}^{+ \infty}
        x \left[ k\right] \cdot y \left[ n - k\right]
    \end{equation}
}

allora, possiamo scrivere che nel dominio di Fourier vale la seguente proprietà di convoluzione: 

{
    \Large 
    \begin{equation}
        x \left[ n\right]
        * 
        y \left[ n\right] 
        = 
        X \left( e^{\jmath \omega}\right) 
        \cdot 
        Y \left( e^{\jmath \omega}\right)  
    \end{equation}
}

\newpage 

\textbf{Modulazione} \newline 

{
    \Large 
    \begin{equation}
        x \left[ n\right]
        \cdot 
        y \left[ n\right] 
        = 
        \dfrac{1}{2 \pi} 
        \int_{-\pi}^{+ \pi}
        X \left( e^{\jmath \theta}\right)
        \cdot 
        Y \left( e^{\jmath \left( \omega - \theta\right)}\right)
        d\theta
    \end{equation}
}

Grazie agli estremi dell'integrale $-\pi$ e $+\pi$, 
sappiamo che, nel dominio di $\omega$, 
la proprietà di modulazione esprime che si sta facendo una convoluzione periodica. \newline 

\textbf{Teorema di Parseval} \label{teo: Teorema di Parseval}\newline 

Il Teorema di Parseval esprime che se: 

{
    \Large 
    \begin{equation}
        \sum_{n = -\infty}^{+ \infty} 
        \abs{x \left[ n\right]}^{2}
        = 
        \dfrac{1}{2 \pi} 
        \int_{-\pi}^{+ \pi} 
        \abs{X \left( e^{\jmath \omega}\right)}^{2}
        d\omega
    \end{equation}
}

dove $\abs{X \left( e^{\jmath \omega}\right)}^{2}$ è lo spettro di densità di energia della sequenza, 
allora vale la seguente operazione: 

{
    \Large 
    \begin{equation}
        \sum_{n = - \infty}^{+ \infty}
        x \left[ n\right]
        \cdot 
        y^{*} \left[n \right]
        = 
        \dfrac{1}{2 \pi}
        \int_{- \pi}^{+ \pi}
        X \left(e^{\jmath \omega}\right)
        \cdot 
        Y^{*} \left(e^{\jmath \omega}\right) 
        d\omega
    \end{equation}
}

\begin{tcolorbox}
Siccome l'energia è qualcosa di fisico, 
formalmente non possiamo scrivere che $\abs{X \left( e^{\jmath \omega}\right)}^{2}$ è lo spettro di energia. \newline 

Ma, per noi, $\abs{X \left( e^{\jmath \omega}\right)}^{2}$ è lo spettro di energia perchè ci viene garantito dal Teorema di Parseval.
\end{tcolorbox}

\newpage 

\subsection{Simmetria della trasformata di Fourier} \label{subsec: Simmetria della trasformata di Fourier}

\footnote{ 
    Slide del prof | 3 - DTFT | pag 11\\
    Appunti | 2025-10-08 | pag 1 \\ 
    Latex | 2025-10-13 Sera | pag 71
} 

Come nel caso continuo, anche nella trasformata di Fourier nel tempo discreto valgono le seguenti proprietà: 

{
    \Large 
    \begin{equation}
        \setlength{\jot}{10pt} % affecting the line spacing in the environment
        \begin{split}
            x^{*} \left[ n \right]
            &\leftrightarrow 
            X^{*} \left( e^{-\jmath \omega}\right)
            \\
            x^{*} \left[ -n \right]
            &\leftrightarrow 
            X^{*} \left( e^{\jmath \omega}\right)
        \end{split}
    \end{equation}
}

con le quali si dimostra che: 

{
    \Large 
    \begin{equation}
        \setlength{\jot}{10pt} % affecting the line spacing in the environment
        \begin{cases}
        \mathbb{R} \left\{ x \left[ n\right]\right\} 
        &\leftrightarrow 
        X_e \left( e^{\jmath \omega}\right)
        \\
        x_e \left[ n\right] 
        &\leftrightarrow 
        \mathbb{R} \left\{ X \left( e^{\jmath \omega}\right)\right\}
        \end{cases}
    \end{equation}
}

e: 

{
    \Large 
    \begin{equation}
        \setlength{\jot}{10pt} % affecting the line spacing in the environment
        \begin{cases}
        \jmath \mathcal{J} \left\{ x \left[ n\right]\right\} 
        &\leftrightarrow 
        X_o \left( e^{\jmath \omega}\right)
        \\
        x_o \left[ n\right] 
        &\leftrightarrow 
        \jmath \mathcal{J} \left\{ X \left( e^{\jmath \omega}\right)\right\}
        \end{cases}
    \end{equation}
}

dove: 

{
    \Large 
    \begin{equation}
        X \left(e^{\jmath \omega }\right)
        = 
        X_e \left(e^{\jmath \omega}\right)
        +
        X_o \left(e^{\jmath \omega}\right) 
    \end{equation}
}

che a sua volta, grazie alle formule di Eulero, possiamo scrivere le componenti: 

{
    \Large 
    \begin{equation}
        \setlength{\jot}{10pt} % affecting the line spacing in the environment
        \begin{cases}
        X_e \left(e^{\jmath \omega}\right)
        &=
        \dfrac{X \left(e^{\jmath \omega}\right) + X^{*} \left(e^{-\jmath \omega}\right)}{2}  
        \\
        X_o \left(e^{\jmath \omega}\right)
        &=
        \dfrac{X \left(e^{\jmath \omega}\right) - X^{*} \left(e^{-\jmath \omega}\right)}{2}    
        \end{cases}
    \end{equation}
}

Per sequenze $x \left[ n\right]$ reali: 

{
    \Large 
    \begin{equation}
        X \left( e^{\jmath \omega}\right)
        = 
        X^{*} \left( e^{-\jmath \omega}\right)
    \end{equation}
}

e che quindi dimostra la simmetria congiunta. \newline 

Siccome $X \left( e^{\jmath \omega}\right)$ è simmetrica, possiamo scrivere che: 

\begin{itemize}
    \item $\mathbb{R} \left\{ X \left( e^{\jmath \omega}\right)\right\}$ e $\abs{X \left( e^{\jmath \omega}\right)}$ sono funzioni pari 
    \item $\mathcal{J} \left\{ X \left( e^{\jmath \omega}\right)\right\}$ e $\sphericalangle X \left( e^{\jmath \omega}\right)$ sono funzioni dispari
\end{itemize}

\newpage 

\begin{tcolorbox}
    Ricordo al volo cosa significa avere una funzione pari e dispari dal punto di vista grafico: 
    \includegraphics[scale = 0.3]{funzione pari e dispari.png}
\end{tcolorbox} 

Inoltre, grazie alla simmetria che godono le sequenze reali, 
possiamo analizzare la trasformata solo nella metà dell'intervallo delle pulsazioni, 
quindi piuttosto che analizzare la trasformata in: 

{
    \Large 
    \begin{equation}
        -\pi \leq \omega \leq + \pi
    \end{equation}
}

possiamo analizzarla nell'intervallo: 

{
    \Large 
    \begin{equation}
        0 \leq \omega \leq +\pi
    \end{equation}
}

\newpage 

\subsection{Alcune trasformate di Fourier} \label{subsec: Alcune trasformate di Fourier}

\footnote{ 
    Slide del prof | 3 - DTFT | pag 12\\
    Appunti | 2025-10-08 | pag 1 
} 

\begin{comment}
\begin{tcolorbox}[colback=green!5!white,colframe=green!75!black]
TO-DO: Fare le foto dei grafici su Matlab e poi metterle in questa subsection 
\end{tcolorbox} 
\end{comment}


Di seguito un elenco delle trasformate di Fourier a tempo discreto che andremo ad utilizzare spesso durante il corso. \newline 

{
    \Large 
    \begin{equation}
        \sum_{n = - \infty}^{+ \infty}
        \delta \left[n\right] 
        \leftrightarrow 
        1
    \end{equation}
}

cioè la trasformata di Fourier di una serie di impulsi unitari nel dominio n 
è uguale a un impulso unitario nel dominio di Fourier. \newline 


{
    \Large 
    \begin{equation}
        \sum_{n = - \infty}^{+ \infty}
        \delta \left[ n - n_0\right]
        \leftrightarrow 
        e^{-\jmath \omega n_0}
    \end{equation}
}

cioè la trasformata di Fourier di una serie di impulsi unitari ritardati di $n_0$ nel dominio n 
è uguale ad un coseno ritardato di $n_0$ nel dominio di Fourier. \newline 

{
    \Large 
    \begin{equation}
        1
        \leftrightarrow
        \sum_{k = -\infty}^{+ \infty}
        2 \pi 
        \delta \left( \omega + 2 \pi k\right)
    \end{equation}
}

dove, con questa notazione, si vuole intendere che un impulso unitario nel dominio n significa in Fourier 
una sommatoria di Delta di Dirac $\delta$ ogni $\omega + 2 \pi k$: in altre parole $\delta \left( \omega + 2 \pi k\right)$ è un treno di impulsi il cui valore tende a $+ \infty$ nelle ampiezze. \newline 

{
    \Large 
    \begin{equation}
        \sum_{n = - \infty}^{+ \infty}
        e^{\jmath \omega_0 n}
        \leftrightarrow 
        \sum_{k = -\infty}^{+ \infty}
        2 \pi 
        \delta \left( \omega - \omega_0 + 2 \pi k\right)
    \end{equation}
}

dove, con questa notazione, si vuole intendere che una serie di sinusoidi a pulsazione fondamentale di $\omega_0$ nel dominio n è uguale a 
una sommatoria di Delta di Dirac $\delta$ ogni $\omega -\omega_0 + 2 \pi k$: in altre parole $\delta \left( \omega - \omega_0 + 2 \pi k\right)$ è un treno di impulsi il cui valore tende a $+ \infty$ nelle ampiezze. \newline 

{
    \Large 
    \begin{equation} \label{eq: trasformata del sinc discreto nel dominio n}
        \sum_{n = - \infty}^{+ \infty}
        \dfrac{\sin \left[ \omega_c n\right]}{\pi n}
        \leftrightarrow 
        \begin{cases}
            \setlength{\jot}{10pt} % affecting the line spacing in the environment
            1 \quad &\text{per } \abs{\omega} < \omega_c 
            \\
            0 \quad &\text{per } \omega_c < \abs{\omega} \leq \pi 
        \end{cases}
    \end{equation}
}



cioè, la trasformata del sinc discreto nel dominio n è una finestra nel dominio di Fourier. \newline 

\newpage 

{
    \Large 
    \begin{equation}
        \setlength{\jot}{10pt} % affecting the line spacing in the environment
        \begin{cases}
            1 \quad \text{per } 0 \leq n \leq M 
            \\
            0 \quad \text{altrove }
        \end{cases} 
        \leftrightarrow
        \dfrac{\sin \left[ \dfrac{\omega \left(M + 1\right)}{2}\right]}{\sin \left[ \dfrac{\omega}{2}\right]}
        e^{\frac{- \jmath \omega M}{2} }
    \end{equation}
}

cioè un gradino (che è una sequenza causale) nel dominio n è "quasi" uguale ad un sinc discreto $\dfrac{\sin \left[ \dfrac{\omega \left(M + 1\right)}{2}\right]}{\sin \left[ \dfrac{\omega}{2}\right]}$ 
ritardato di $e^{\frac{- \jmath \omega M}{2} }$. \newline 

Ultima funzione, ma non per meno importanza: 

{
    \Large 
    \begin{equation}
    \sum_{n = - \infty}^{+\infty}
    a^{n} \cdot u\left[ n\right] \text{ con } \left( \abs{a < 1}\right)
    \leftrightarrow 
    \dfrac{1}{1 - a \cdot e^{- \jmath \omega}}         
    \end{equation}
}

cioè un esponenziale $a^n$ moltiplicato per un gradino $u \left[ n\right]$ nel dominio n 
è uguale a una funzione razionale reale. \newline 

\begin{tcolorbox}
    Le funzioni razionali reali le abbiamo viste ad Elettrotecnica quando si discuteva della trasformata di Laplace. \newline 

    Per un breve ripasso della trasformata di Laplace: \newline 

    But what is a Laplace Transform? by 3Blue1Brown \\
    \url{https://www.youtube.com/watch?v=j0wJBEZdwLs&t=452s}

\end{tcolorbox}

\newpage 

\section{Filtri ideali}

\footnote{ 
    Slide del prof | 3 - DTFT | pag 15 - 18\\
    Appunti | 2025-10-08 | pag 1 - 2
} 

Essendo "nel mondo ideale del DSP", 
possiamo sagomare un sistema discreto che permette di sagomare lo spettro del segnale di ingresso. \newline 

Prima di introdurre i filtri ideali, 
vorrei ricordare (per almeno la decima volta) che abbiamo a che fare con sistemi LTI, 
cioè sistemi lineari a tempo invariante (Linear Time-Invariant). \newline  

Essendo sistemi lineari, quest'ultimi possono solo togliere delle componenti in frequenza e non aggiungerne di nuovi (per quello servono i sistemi non lineari). \newline 

Inoltre, in questo corso, 
andremo ad analizzare sequenze reali che hanno spettro periodico che si estende in: 

{
    \Large 
    \begin{equation}
        0 \leq \omega \leq + \pi
    \end{equation}
}

Dopo queste introduzioni, possiamo introdurre i primi filtri ideali. \newline 

\textbf{Filtro passa-basso (Low Pass Filter)} \newline 

La risposta in frequenza di un filtro passa basso $H_{lp} $ ha la seguente funzione: 

{
    \Large 
    \begin{equation}
        H_{lp} \left(e^{\jmath \omega}\right) 
        = 
        \begin{cases}
            1 \text{ per } \abs{\omega} < \omega_c 
            \\
            0 \text{ per } \omega_c < \abs{\omega} \leq \pi
        \end{cases}
    \end{equation}
}

dove $\omega_c$ è la pulsazione di taglio. \newline 

Da un punto di vista grafico, lo spettro in $\omega$ di un filtro passa basso è il seguente: 

\begin{figure}[h]
    \centering
    \includegraphics[scale = 0.4]{Spettro di un filtro passa basso ideale.png}
\end{figure}  

Come si nota dallo spettro (colorato di blu in figura), 
lo spettro di un filtro passa-basso è periodico (si ripete) ogni $k \pi$. \newline

\newpage 

\textbf{Filtro passa alto (High Pass Filter)} \newline 

La risposta in frequenza di un filtro passa alto $H_{hp} $ ha la seguente funzione: 

{
    \Large 
    \begin{equation}
        H_{hp} \left(e^{\jmath \omega}\right) 
        = 
        \begin{cases}
            1 \text{ per } \omega_c < \abs{\omega} \leq \pi
            \\
            0 \text{ per } \abs{\omega} < \omega_c 
        \end{cases}
    \end{equation}
}

dove $\omega_c$ è la pulsazione di taglio. \newline 

Da un punto di vista grafico, lo spettro in $\omega$ di un filtro passa alto è il seguente: 

\begin{figure}[h]
    \centering
    \includegraphics[scale = 0.4]{Spettro di un filtro passa alto ideale.png}
\end{figure}  

Come si nota dallo spettro (colorato di blu in figura), 
lo spettro di un filtro passa-alto è periodico (si ripete) ogni $k \pi$. \newline

\textbf{Filtro passa banda (Band Pass Filter)} \newline 

La risposta in frequenza di un passa banda alto $H_{bp} $ ha la seguente funzione: 

{
    \Large 
    \begin{equation}
        H_{bp} \left(e^{\jmath \omega}\right) 
        = 
        \begin{cases}
            1 \text{ per } \omega_a < \abs{\omega} \leq \omega_b
            \\
            0 \text{ altrove tra } -\pi, +\pi 
        \end{cases}
    \end{equation}
}

dove $\omega_a$ e $\omega_b$ sono le pulsazioni che il filtro permette di far passare. \newline 

Da un punto di vista grafico, lo spettro in $\omega$ di un filtro passa basso è il seguente: 

\begin{figure}[h]
    \centering
    \includegraphics[scale = 0.4]{Spettro di un filtro passa banda ideale.png}
\end{figure}  

Come si nota dallo spettro (colorato di blu in figura), 
lo spettro di un filtro passa-banda è periodico (si ripete) ogni $k \pi$. \newline

\newpage 
\textbf{Elimina banda (Stop band Filter)} \newline 

La risposta in frequenza di un filtro elimina banda $H_{bs} $ ha la seguente funzione: 

{
    \Large 
    \begin{equation}
        H_{bs} \left(e^{\jmath \omega}\right) 
        = 
        \begin{cases}
            1 \text{ altrove tra } -\pi, +\pi 
            \\
            0  \text{ per } \omega_a < \abs{\omega} \leq \omega_b
        \end{cases}
    \end{equation}
}

dove $\omega_a$ e $\omega_b$ sono le pulsazioni che il filtro non permette di far passare. \newline 

Da un punto di vista grafico, lo spettro in $\omega$ di un filtro passa basso è il seguente: 

\begin{figure}[h]
    \centering
    \includegraphics[scale = 0.4]{Spettro di un filtro elimina banda ideale.png}
\end{figure}  

Come si nota dallo spettro (colorato di blu in figura), 
lo spettro di un filtro elimina-banda è periodico (si ripete) ogni $k \pi$. \newline

\newpage 

\begin{comment}
\subsection{Trasformata di Fourier in Matlab}

\footnote{ 
    Slide del prof | 3 - DTFT | pag 13 - 14\\
    Appunti | 2025-10-08 | pag 3 
} 

\begin{tcolorbox}[colback=green!5!white,colframe=green!75!black]
TO-DO: Spiegazione codice Esempio 3.1, 3.2 
\end{tcolorbox} 

\newpage 

\end{comment}

\subsection{Filtri passa basso ideale} \label{subsec: Filtri passa basso ideale}

\footnote{ 
    Slide del prof | 3 - DTFT | pag 19\\
    Appunti | 2025-10-08 | pag 2 - 3
} 

Come abbiamo visto dallo spettro dei filtri passa alto, passa banda ed elimina banda, 
tutti questi filtri possono essere ottenuti da un filtro passa-basso (opportunamente shiftato nello spettro). \newline 

Ricordando l'equazione \ref{eq: trasformata del sinc discreto nel dominio n} a pag \pageref{eq: trasformata del sinc discreto nel dominio n},  
di seguito scriviamo la trasformata di una finestra nel dominio di $\omega$ è:  

{
    \Large 
    \begin{equation}
        H_{lp} \left( e^{\jmath \omega}\right)
        =
        \begin{cases}
            1 \text{ per } \abs{\omega} < \omega_c 
            \\
            0 \text{ per } \omega_c < \abs{\omega} \leq \pi 
        \end{cases}
        \leftrightarrow 
        h_{lp} \left[ n\right]
        =
        \sum_{n = - \infty}^{+ \infty}
        \dfrac{\sin \left[ \omega_c n\right]}{\pi n}
    \end{equation}
}

cioè è un sinc discreto che si estende in nel dominio n nell'intervallo $\left[ - \infty, + \infty\right]$. \newline 

Possiamo rappresentare graficamente il sinc discreto con la seguente figura: 

\begin{figure}[h]
    \centering
    \includegraphics[scale = 0.5]{Sinc discreto.png}
\end{figure}  

ricordando che la sequenza si estende tra $- \infty$ e $+ \infty$. \newline 

Il fatto che il filtro passa basso (cioè lo spettro di una finestra ideale nello spettro), nel dominio delle sequenze n, 
si estende tra $\left[ - \infty, + \infty\right]$, rende il filtro passa-basso un sistemi LTI IIR (Infinite Impulse Reponse) e non causale: 
questo è un problema perchè, teoricamente parlando, se dovessimo calcolare una convoluzione tra due sequenze su Matlab, ci impiegherebbe un tempo infinito. \newline 

Ma, ponendo che: 

{
    \Large 
    \begin{equation}
        \lim_{M \to +\infty} 
        \int_{- \pi}^{+ \pi} 
        \abs{
            H_{lp} \left( e^{\jmath \omega}\right)
            - 
            \sum_{n = - M}^{+ M}
            \dfrac{\sin \left[ \omega_c n\right]}{\pi n}
            e^{- \jmath \omega n}
        }^{2}
        d \omega 
        = 
        0
    \end{equation}
}

sappiamo che per un M sufficientemente alto, avremo che il sinc discreto nel dominio n tende a una finestra nel dominio delle frequenze. \newline 

In particolare, la funzione appena scritta esprime che esiste una convergenza in senso quadratico tra 
la differenza tra un filtro passa-basso ideale nel dominio di $\omega$ e la trasformata di un sinc discreto per M che tende ad infinito: questo fenomeno prende il nome di fenomeno di Gibbs. \newline 

\newpage 

Da un punto di vista grafico vedremo il seguente caso: 

\begin{figure}[h]
    \centering
    \includegraphics[scale = 0.5]{Fenomeno di Gibbs rispetto ad un filtro passa basso ideale.png}
\end{figure}  

dove in blu è tratteggiato l'andamento dello spettro in $\omega$ di un filtro passa basso ideale ed in rosso lo spettro di un filtro passa basso non ideale con M molto grande. \newline 

\begin{tcolorbox}
    Per visualizzare meglio cosa avviene realmente per un M sufficiente alto, 
    ti consiglio di visualizzare questo video (con anche un esempio su Matlab): \newline
    
    Fourier Series and Gibbs Phenomena [Matlab] by Steve Brunton\\
    \url{https://www.youtube.com/watch?v=98VixA3MjAc}
\end{tcolorbox}

\newpage 



